\documentclass[11pt,a4paper]{article}
\newcommand{\CvKind}{Research CV}
% Shared layout for the digital and print editions. Run from the repository root.
\usepackage[a4paper,margin=18mm,footskip=10mm]{geometry}
\usepackage{fontspec}
\setmainfont{XCharter}[
  Extension=.otf,
  UprightFont=*-Roman,
  BoldFont=*-Bold,
  ItalicFont=*-Italic,
  BoldItalicFont=*-BoldItalic,
  RawFeature={-kern,-liga,-clig},
  Ligatures=TeX
]
\usepackage{titlesec,enumitem,ragged2e,xcolor,fancyhdr,lastpage}
\definecolor{Accent}{HTML}{2446B3}
\definecolor{Muted}{HTML}{475569}
\ifdefined\PrintEdition
  \definecolor{Accent}{gray}{0.12}
  \definecolor{Muted}{gray}{0.25}
  \usepackage{qrcode,tikz}
\fi
\usepackage[pdfauthor={Le Nguyen Gia Hung},colorlinks=true,urlcolor=Accent,linkcolor=Accent]{hyperref}
\hypersetup{pdftitle={Le Nguyen Gia Hung - \CvKind}}
\pagestyle{fancy}
\fancyhf{}
\renewcommand{\headrulewidth}{0pt}
\renewcommand{\footrulewidth}{0pt}
\fancyfoot[L]{%
  \ifdefined\PrintEdition
    \raisebox{-1.8mm}{\input{cv-running-duck.tex}}\hspace{2mm}%
  \fi
  \fontsize{9}{11}\selectfont\textcolor{Muted}{Le Nguyen Gia Hung\enspace|\enspace\CvKind}}
\fancyfoot[R]{\fontsize{9}{11}\selectfont\textcolor{Muted}{\thepage\,/\,\pageref*{LastPage}}}
\setcounter{secnumdepth}{0}
\setlength{\parindent}{0pt}
\raggedright
\setlength{\emergencystretch}{2em}
\hyphenpenalty=10000
\exhyphenpenalty=10000
\titleformat{\section}{\fontsize{13}{15}\selectfont\bfseries\color{Accent}}{}{0pt}{}[{\vspace{1pt}\titlerule[0.35pt]}]
\titlespacing{\section}{0pt}{6pt}{4pt}
\newenvironment{highlights}{\begin{itemize}[topsep=2pt,parsep=0pt,partopsep=0pt,itemsep=0pt,leftmargin=12pt]}{\end{itemize}}
\renewcommand\labelitemi{\textbullet}
\newcommand{\resumeEntry}[4]{%
  \par\noindent\textbf{#1}\hfill{\small\textcolor{Muted}{#4}}\par
  {\small\textit{#3}\hfill\textcolor{Muted}{#2}}\par\vspace{1pt}}
\newcommand{\paperGap}{\par\vspace{3pt}}
\newcommand{\cvHeader}[2]{%
  \ifdefined\PrintEdition\begin{minipage}[t]{0.83\linewidth}\vspace{0pt}\fi
  {\fontsize{25}{28}\selectfont\bfseries\textcolor{Accent}{Le Nguyen Gia Hung}}\par\vspace{3pt}
  {\fontsize{11}{13}\selectfont #1}\par\vspace{2pt}
  {\fontsize{10}{12}\selectfont #2}\par\vspace{5pt}
  {\fontsize{10}{12}\selectfont
    \href{mailto:heiontheway@gmail.com}{heiontheway@gmail.com}\enspace|\enspace
    \href{tel:+84394262651}{+84 394 262 651}\enspace|\enspace Ho Chi Minh City, Vietnam\par
    \href{https://github.com/hei1sme}{github.com/hei1sme}\enspace|\enspace
    \href{https://hei-portfolio.vercel.app}{hei-portfolio.vercel.app}\enspace|\enspace
    \href{https://linkedin.com/in/le-nguyen-gia-hung}{linkedin.com/in/le-nguyen-gia-hung}\par
  }
  \ifdefined\PrintEdition
    \end{minipage}\hfill\begin{minipage}[t]{0.13\linewidth}\vspace{0pt}\centering
    \qrcode[height=18mm,level=M]{https://hei-portfolio.vercel.app}\par\vspace{3pt}
    {\fontsize{9}{11}\selectfont Portfolio}
    \end{minipage}\par
  \fi
  \vspace{4pt}
}

\begin{document}
\cvHeader{AI undergraduate researcher\enspace|\enspace Explainable AI \& multimodal learning}
{Research interests: trustworthy AI, time-series forecasting and computer vision.}
{\fontsize{10}{12}\selectfont
\href{https://orcid.org/0009-0003-7120-8167}{ORCID: 0009-0003-7120-8167}\enspace|\enspace
\href{https://scholar.google.com/citations?user=iGODAQYAAAAJ}{Google Scholar: iGODAQYAAAAJ}\par}

\section{Education}
\textbf{FPT University} --- B.Eng. in Artificial Intelligence\hfill{\small 2023 -- Dec 2027 (expected)}\par
GPA: \textbf{8.62/10.0}. Certificate of Merit: semesters 3, 4, 5 and 6.

\section{Publications \& Manuscripts}
\textbf{Peer-reviewed conference publication}\par
\textbf{Gia-Hung Nguyen Le}, Gia-Bao Pham Hoang, Tan-Phat Vo, Thu Le and Nhu Nguyen.
``Proactive Air Quality Forecasting and Health Alert System for Melbourne.''
\textit{AJCAI 2025, Springer LNAI}. First and corresponding author.
\href{https://doi.org/10.1007/978-981-95-4969-6_34}{DOI: 10.1007/978-981-95-4969-6\_34}.
\paperGap
\textbf{Manuscripts under review}\par
\textbf{Gia-Hung Nguyen Le}, Thu Le and Huynh Quoc Viet.
``Distribution-Free Coverage of Attack Classes on Devices Absent from Training.''
\textit{AISTATS 2027, under review}. First author.
\paperGap
\textbf{Gia-Hung Nguyen Le}, Quoc-Viet Huynh and Thu Le.
``Auditing Multi-Source Cyanobacterial Bloom Forecasts under Sparse Environmental Telemetry.''
\textit{ACIIDS 2027, under review}. First and corresponding author.
\paperGap
Thu Le, \textbf{Le Nguyen Gia Hung}, Huynh Quoc Viet, Nguyen Hoai My and Trieu Vuong Nguyen.
``Trustworthy Multimodal Anomaly Detection in Time Series with Conformal Calibration and Counterfactual Attribution.''
\textit{ACM SIGKDD 2027, under review}. Second author.
\paperGap
Thu Le, \textbf{Gia-Hung Nguyen Le}, Hoang Pham Gia Bao and Huynh Quoc Viet.
``Warning Before the Evening Peak: Health-Centred Air Quality Alerts in New Zealand and Australia.''
\textit{AJCAI 2026, under review}. Second author.

\section{Selected Research Experience}
\resumeEntry{SmokeNet: Air Quality Forecasting}{Academic research}
{First \& Corresponding Author; published at AJCAI 2025}{Jun -- Nov 2025}
\begin{highlights}
\item Designed a Transformer using 48 hours of history to forecast next-day mean PM\textsubscript{2.5}, with seasonal embeddings and reconstruction loss.
\item Reported MAE \textbf{0.7470} $\mu$g/m$^3$, $R^2$ \textbf{0.9545} and \textbf{35.9\% lower MAE than XGBoost} on the held-out 2024 test year.
\item Examined forecasts with SHAP, Integrated Gradients and DiCE; a separate May 2024 smoke-event backtest reported \textbf{57.7\% lower MAE}. \href{https://github.com/SpeedyLabX/melbourne-air-quality-forecast}{Code}.
\end{highlights}

\resumeEntry{Cyanobacterial Bloom Forecast Audit}{Academic research}
{First \& Corresponding Author; ACIIDS 2027 under review}{Jan 2026 -- Present}
\begin{highlights}
\item Co-authored a past-only audit across \textbf{7 New Zealand lakes}, evaluating 7- and 14-day forecasts on \textbf{170 observed holdout targets} per horizon.
\item Full fusion showed no significant gain over Chl-a history after Holm correction; provenance checks reduced 22 apparent episode starts to 3 eligible incidents.
\end{highlights}

\resumeEntry{Device-Held-Out Conformal Prediction}{Academic research}
{First Author; AISTATS 2027 under review}{}
\begin{highlights}
\item Co-authored a device-disjoint conformal prediction study on CIC IoT-DIAD 2024: \textbf{40 training, 14 calibration and 14 test devices}; compared device-weighted and device-quantile rules. \href{https://github.com/SpeedyLabX/iot-diad-device-held-out}{Code}.
\end{highlights}

\newpage
{\small\textcolor{Muted}{Research experience continued}}\par\vspace{5pt}

\resumeEntry{Trustworthy Multimodal Anomaly Detection}{Academic research}
{Second Author; ACM SIGKDD 2027 under review}{Jan 2026 -- Present}
\begin{highlights}
\item Co-authored a leakage audit across \textbf{9 domains} and the C-DLE framework for calibrated false-alarm control; reported \textbf{94.3\% attribution accuracy} against ground-truth anomalies.
\end{highlights}

\textbf{Other research.} Second author on health-centred air-quality alerts: the study reported 1-day-ahead MAE 3.19 $\mu$g/m$^3$ and $R^2=0.54$ in Canterbury.\par
\resumeEntry{MERR-GAT: Multimodal Emotion Recognition}{Research in progress}
{Second \& Corresponding Author}{Oct 2025 -- Present}
\begin{highlights}
\item Develop audio--text emotion recognition with \textbf{WavLM + Compact Conformer}, \textbf{BERT/RoBERTa} and GAT fusion; evaluation on IEMOCAP, RAVDESS and MELD is ongoing.
\end{highlights}
\resumeEntry{Vietnamese Speech Recognition}{Research in progress}
{Co-Researcher \& Co-Author}{Sep 2025 -- Present}
\begin{highlights}
\item Develop a \textbf{5.4M-parameter} CNN--BiLSTM/CTC model on \textbf{745 hours} of speech; reported \textbf{WER 33.09\%} and \textbf{CER 15.28\%}.
\end{highlights}
\textbf{Selected course project:} Built 9-class PPE detection with YOLOv8/OpenCV, compliance logic and a Streamlit webcam interface with violation reports.

\section{Experience \& Leadership}
\resumeEntry{First Cloud AI Journey --- NutriTrack}{AWS Vietnam program}
{FCAJ Cloud Intern, Team NeuraX}{Jan -- May 2026}
\begin{highlights}
\item Developed recipe-generation and AI-coaching workflows using \textbf{Amazon Bedrock and AWS Lambda}, structured JSON responses and DynamoDB user context.
\item Collaborated on AppSync integration; documented timeout handling and workshops. \href{https://hei1sme.github.io/hei-FCAJ-intership-report/}{Report}.
\end{highlights}
\resumeEntry{SpeedyLabX Student Research Group}{FPT University}
{Founder \& Student Lead}{Jan 2025 -- Present}
\begin{highlights}
\item Lead a \textbf{10-member} undergraduate group through technical workshops, paper reading and shared research and submission workflows.
\end{highlights}

\section{Hackathon Projects \& Competitions}
\resumeEntry{BioListen VN --- Vietnam AI Innovation Challenge}{Team NeuraX.ai}
{UI \& LLM; placed 6th in Round 1}{Jul 2026}
\begin{highlights}
\item Built monitoring views for audio input, spectrograms and alerts; integrated API data and Vietnamese ranger-report prompts. Next.js, TypeScript, FastAPI/Groq stack. \href{https://github.com/hei1sme/BioListen-VN}{Code}.
\end{highlights}
\resumeEntry{CosmicOptic --- NASA Space Apps Challenge}{Hackathon participant}
{Frontend \& Backend Developer}{Oct 2025}
\begin{highlights}
\item Built React/D3 light-curve views and a FastAPI backend; demonstrated the inference interface with synthetic signals. \href{https://github.com/hei1sme/cosmicoptic-webapp}{Code}.
\end{highlights}
\textbf{Also participated:} FPT Hackathon 2025 (algorithmic programming); FPT University AI Hackathon (TeamFinder, a course-teammate website).

\section{Skills}
\textbf{Languages \& frameworks:} Python, SQL/PostgreSQL, TypeScript, React, FastAPI.\\
\textbf{ML \& evaluation:} PyTorch, Transformers, TensorFlow/Keras, Scikit-learn, OpenCV; SHAP, Integrated Gradients, DiCE, chronological evaluation.\\
\textbf{Data \& tools:} Pandas, NumPy, Git, Docker, Linux; Bedrock, Lambda, DynamoDB, AppSync.

\section{Selected Certifications}
Natural Language Processing (DeepLearning.AI); Gradient to Production: MLOps \& Model Serving (Coursera); Full Stack Software Developer (IBM).
\end{document}
